\section{Scalar alternating and symmetric inducing blocks}
\begin{theorem}[D-BI-08: scalar standard inducing representations]
\label{D-BI-08}
Let an exact-$K$ Weil sixfold representation be induced from a
complex representation $\tau$ whose finite image, in its given
representation, satisfies
\[
 \mathfrak A_n\subset J\subset\mathbf C^*\mathfrak S_n,
 \qquad \tau=\chi_d\otimes(\operatorname{Std}_n\circ\pi),
 \qquad n\ge4,
\tag{BI.12}
\]
where the permutation groups act on their deleted permutation space
and $\pi:J\to\mathfrak S_n$ and $\chi_d:J\to\mu_d$ are the
permutation and scalar projections, with scalar image of order $d$.
A rank-nine period
map forces $d\in\{3,4,6\}$, $K=\mathbf Q(\zeta_d)$, and split
actual polarization.  In particular this cannot supply a dominant
nonsplit family.  There is no assumption that the inducing
cyclotomic field initially equals $K$.

This applies, in particular, to an inducing representation whose
reflection normal closure is an irreducible group of type $A_{n-1}$:
its complex linear normalizer is $\mathbf C^*\mathfrak S_n$.
It also includes scalar extensions of the alternating deleted
representation in the Guralnick--Saxl bireflection classification;
a direct-product assumption is not being imposed.
\end{theorem}
\begin{proof}
Set $k=n-1$ and $L=\mathbf Q(\zeta_d)$.  The displayed matrices
give a model of $\tau$ over $L$; induced characters give
$K\subset\mathbf Q(\chi_\tau)\subset L$, so $d>2$.
The scalar and permutation projections are well-defined because
the deleted permutation group contains no nonidentity scalar.
Quotient the kernel of $\tau$ in its
stabilizer to obtain a connected $J$-cover $C_0\to X_0$.
Frobenius reciprocity identifies its $\tau$-multiplicity variation
with the original complex sixfold multiplicity, exactly as in
D-BI-01.  If $h=g(X_0)$, Euler characteristic gives
$D=6-2k(h-1)$, where $D$ is the total moving dimension.
For $h=1$ the marked-base moduli dimension is at most six.  For
$h\ge2$ and $k\ge3$ the only nonnegative case is $h=2,k=3,D=0$,
whose base moduli has dimension three.  Therefore rank nine forces
$h=0$, at least twelve genuine branch points, and
\[
       \sum_j\bigl(\operatorname{rank}(1-g_j)-1\bigr)
                    =2k+6-r\le2n-8.
\tag{BI.13}
\]

Write $g_j=\lambda_j\pi_j$ on the deleted permutation space.
If the order of $\lambda_j\ne1$ is $e_j$, its fixed space has
dimension at most $\lfloor n/e_j\rfloor$: only permutation cycles
whose lengths are divisible by $e_j$ contribute.  Hence its
contribution to (BI.13) is at least
$n-2-\lfloor n/e_j\rfloor$.
Let $t$ count the nonreal scalar entries and $u$ the entries
with scalar $-1$.  There must be a nonreal entry, since these
entries generate the nonreal scalar quotient; their scalar
product is one, so $t\ne1$.

If $t\ge3$, three of them alone exceed (BI.13), since
\[
 3(n-2-\lfloor n/3\rfloor)-(2n-8)
                   =n+2-3\lfloor n/3\rfloor\ge2.
\]
Thus $t=2$.  If $u\ge2$, two nonreal entries and two minus-one
entries exceed (BI.13), with positive difference
\[
               2n-2\lfloor n/3\rfloor-2\lfloor n/2\rfloor>0.
\]
If $u=1$, the two nonreal scalars have product $-1$.
Their orders $e_1,e_2\ge3$ satisfy
$1/e_1+1/e_2\le1/2$: if one has order three the other has order
six, and otherwise both orders are at least four.  The three
entries therefore contribute at least $2n-6$, again larger than
$2n-8$.  Consequently
\[
             t=2,\qquad u=0,\qquad
             \lambda_2=\lambda_1^{-1}.
\tag{BI.14}
\]
Since these scalars generate $\mu_d$, both have order $d$.
The cyclic quotient $T=C_0/\ker\chi_d\to X_0=\mathbf P^1$
is branched at exactly two points, fully ramified at each, and
has genus zero.

Let $H_1$ be the stabilizer of one letter under the permutation
projection of $J$.  Its scalar image is all of $\mu_d$:
for any element of $J$, an element of the real $\mathfrak A_n$
subgroup can move its image of that letter back to the letter,
without changing the scalar.  Put
$Y=C_0/\ker(\chi_d|_{H_1})$ and $X=C_0/H_1$.  Then $Y\to X$
is a connected cyclic cover of degree $d$.  The permutation
representation on the $n$ letters is $1\oplus\operatorname{Std}_n$,
so the representation identity
\[
 \operatorname{Ind}_{H_1}^{J}(\chi_d|_{H_1})
       =\chi_d\oplus\tau
\tag{BI.15}
\]
identifies its faithful-character cohomology with the direct sum
of the scalar-character and $\tau$ multiplicities.  The former
is zero because it comes from $H^1(T)$.

We spell out the rational multiplicity issue, since the inducing
field may be larger than $K$.  For every embedding
$\gamma:L\hookrightarrow\mathbf C$, (BI.15) holds with
$\chi_d^\gamma$ and $\tau^\gamma$, and its scalar-character
summand still has zero cohomology.  Induction from the original
block stabilizer to $G$ sends $\tau^\gamma$ to the original
$\rho$ or its complex conjugate according as $\gamma|_K$ is the
chosen embedding or its conjugate.  There are $\ell=[L:K]$
embeddings of each type.  Thus rational semisimplicity and
Frobenius reciprocity identify the primitive cyclic Hodge structure
with $\ell$ copies of the original primitive factor:
\[
                     P_d(Y/X)\sim B^{\ell}.
\tag{BI.16}
\]
More explicitly, inducing the rational $L$-module underlying
$\chi_d$ gives the sum of $\ell$ copies of $V_{\mathbf Q}$ and
an additional rational representation whose Hom into $H^1(C)$
vanishes.  Taking this Hom discards exactly that zero summand.
The constructions are rational group-algebra correspondences in
families, so (BI.16) is an isogeny of variations, not just an
identity of complex dimensions.
This argument does not require $\mathbf Q(\chi_\tau)=L$:
the multiplicity $\ell$ counts the embeddings of the scalar-character
module, and each induced $\tau^\gamma$ is determined by $\gamma|_K$.

The faithful-character cohomology of $Y\to X$ has dimension six.
Writing $a=g(X)$ and $R$ for its genuine cyclic branch count gives
$6=2a-2+R$.  Its cyclic Hurwitz moduli has dimension at most
$5,6,7,8,9$ for $a=0,1,2,3,4$, respectively.  Equation (BI.16)
identifies its period differential kernel with that of $B$;
therefore rank nine forces $a=4,R=0$, and a dominant image in a
component of the etale cyclic genus-four moduli.  D-MON-02, from
Looijenga's published monodromy theorem and the explicit commutant
argument there, gives $\operatorname{End}^0(P_d)=L$ at a very
general point in that image.  If $\ell>1$, (BI.16) embeds the
noncommutative algebra $M_\ell(K)$ into this field, a contradiction.
Hence $L=K$ and $d=3,4,6$.

The splitness is that of the actual polarization.  Once $L=K$,
the rational faithful-character projector for $H_1$ is
Rosati-fixed.  The character $\chi_d|_{H_1}$ has multiplicity one
in $\tau|_{H_1}$: the subgroup $\mathfrak A_{n-1}$ is transitive
on the other $n-1$ letters, also when $n=4$, so the fixed subspace
of the restricted permutation representation has dimension two,
and that of its deleted summand has dimension one.
By (BI.15) its active image is a single primitive
$\tau$-line: the other line is precisely the zero scalar-character
summand.  After induction the same assertion holds for the original
$\rho$-factor, since any other summand has zero cohomology.
Its image is the pullback of $P_d(Y/X)$, with the Jacobian
polarization scaled by the degree of the quotient map.  Looijenga's
equivariant hyperbolic basis makes the etale cyclic primitive
polarization split.  As in D-MON-03, changing the Rosati-fixed
primitive line only multiplies the multiplicity Hermitian form
by a positive rational scalar and preserves its isotropic
three-space.  This proves the asserted splitness.
\end{proof}

